\documentclass[10pt,a4paper]{amsart}
\usepackage[margin=19mm]{geometry}
\usepackage{amsmath,amssymb,mathtools}
\usepackage[scr=rsfs,cal=euler]{mathalfa}
\usepackage{xfrac}
\usepackage[unicode,hidelinks,bookmarks=false]{hyperref}
\newcommand{\ie}{i.e.\ }
\newcommand{\cf}{cf.\ }
\newcommand{\resp}{respectively\ }
\newcommand{\Subsection}{\subsection}
\providecommand{\nobreakdash}{\nobreak\hbox{-}}
\setlength{\emergencystretch}{2em}
\multlinegap0pt
\usepackage[all]{xy}
\usepackage{enumerate}
\usepackage{bm}
\usepackage{amssymb}
\usepackage{tikz}
\usepackage{mathtools}
\newtheorem{lem}{Lemma}
\DeclareMathOperator{\Li}{Li}
\def\N{\mathbb{N}}
\def\Z{\mathbb{Z}}
\title{\bf Contour Integration and Cyclotomic Ap\'ery-Like Series Involving Generalized Binomial Coefficients}
\author{Ce Xu}
\date{Source: arXiv:2512.18121v2 (CC BY 4.0)}
\begin{document}
\maketitle

\noindent\emph{Source and licence.} \bf Contour Integration and Cyclotomic Ap\'ery-Like Series Involving Generalized Binomial Coefficients. Version arXiv:2512.18121v2 (CC BY 4.0), distributed by arXiv under the Creative Commons Attribution 4.0 International licence, http://creativecommons.org/licenses/by/4.0/, as recorded in the arXiv licence metadata for this exact version and shown on the abstract page https://arxiv.org/abs/2512.18121v2. Full licence notice: \texttt{licenses/arxiv-2512.18121-CC-BY-4.0.txt}.\par\smallskip

\noindent\textit{Editorial scope.} Counted unit: the article's lem lem-extend-xuzhou-two (source0.tex lines 553--557). The article's complete original proof of it is delivered with it (source0.tex lines 558--560, one \texttt{proof} environment of 3 lines). No mathematics is written, repaired or re-derived: the statement and the proof are the article's own lines, in the article's own notation, and no retained line is substituted or re-flowed.  No further source-local context is needed: every cross-reference of every retained line resolves inside the two delivered spans.  Nothing was omitted from any retained span, so the package carries no omission record.  No retained line contains a citation, so the wrapper carries no bibliography and no bibliography entry.  No retained line contains a figure environment, an image-inclusion command or drawing code, so no image is delivered and the package declares no asset dependency.  Editorial formatting only, in addition: an AMS \texttt{amsart} preamble that restates the article's own declarations and macros used by the retained lines, namely the environments \texttt{lem} on separate counters, together with the standard packages those lines need.  The retained environments print in this document as Lemma 1 in order of appearance, whereas the article prints the same items with the numbering of its own full document.  The complete native source of the article as distributed in the arXiv e-print for this exact version is bundled unchanged as \texttt{sources/191362/source0.tex}.
\begin{lem}\label{lem-extend-xuzhou-two} Let $x$ be root of unity. For $n\in\N_0$ and $a\in \mathbb{C}\setminus \Z$, if $|s+n|<1$, then
\begin{align}\label{equ-extend-xuzhou-two}
\Phi(s-a;x)=x^n\sum_{m=0}^\infty \left((-1)^m \Li_{m+1}(x;1-a)-x \Li_{m+1}\Big(x^{-1};a\Big) \right)(s+n)^m.
\end{align}
\end{lem}

\begin{proof}
The proof of \eqref{equ-extend-xuzhou-two} is straightforward from the definition of $\Phi(s;x)$ and is therefore omitted.
\end{proof}

\end{document}
